\section{Zephyr, Tulu 2, SteerLM y Starling: no hay un ganador único}

La sección anterior presentó las fórmulas; esta sección examina la evidencia de los modelos. Zephyr demostró que DPO puede producir un modelo abierto que trasciende el ámbito especializado, y Tulu 2 extendió DPO a 70B; SteerLM y Starling, por su parte, demostraron que la vía de PPO/RL todavía puede ser más fuerte. La conclusión histórica no es «DPO sustituye a PPO», sino que la preference optimization empezó a producir releases públicas comparables.

\subsection{Zephyr y Tulu 2: DPO pasa del artículo a la release}

Este apartado revisa primero dos implementaciones decisivas. Zephyr combina UltraChat y UltraFeedback y le da a DPO una recipe clara y reproducible; Tulu 2 usa los datos y la infraestructura de entrenamiento de AI2 para extender el método a 70B. La publicación de los modelos permite a la comunidad comprobar si el método conserva su valor en distintas escalas y con distintos datos.

\lecturefigure{slide-067.jpg}{Zephyr $\beta$: uno de los primeros modelos abiertos que recibió amplia atención gracias a DPO}{official deck, p.~67}

\lecturefigure{slide-068.jpg}{Tulu 2: extensión de DPO a la escala de 70B}{official deck, p.~68}

\begin{knowledgebox}{Lectura de la figura: las partes transferibles de la recipe}
Al comparar releases hay que descomponerlas en base model, IFT data, preference data, reference policy, $\beta$, sequence length y eval. Si Zephyr y Tulu cambiaron varios componentes a la vez, no se puede atribuir toda la mejora a DPO; su aportación consiste en ofrecer un artifact completo para que los equipos posteriores puedan hacer ablaciones componente por componente.
\end{knowledgebox}

\teachervoice{00:46:33--00:48:30, Nathan describe a Zephyr como el primer modelo con el que DPO realmente «make a splash», y considera a Tulu 2 un hito de escalamiento. La credibilidad de un método ante la comunidad proviene de la release de modelos, no solo de la derivación en un artículo.}

\subsection{SteerLM y Starling: PPO/RL sigue siendo competitivo}

Que DPO sea sencillo no significa que el RL on-policy haya dejado de funcionar. SteerLM controla la salida mediante condicionamiento por atributos y señales de preferencia, y Starling ejecuta PPO sobre datos de preferencia al estilo de Arena, lo que muestra que un pipeline más complejo todavía puede obtener ventajas. El curso coloca deliberadamente estos modelos después del éxito de DPO para impedir el relato de que «un método nuevo lo unifica todo».

\lecturefigure{slide-069.jpg}{SteerLM y Starling: los métodos PPO/RL todavía pueden superar a los modelos DPO}{official deck, p.~69}

\begin{warningbox}{Las clasificaciones de algoritmos suelen disfrazar una ventaja de datos como ventaja del optimizador}
Si el modelo PPO usa preference data más reciente, más grande o más realista, mientras que la baseline DPO usa datos antiguos, el resultado no se puede atribuir al optimizer. Una comparación confiable requiere compartir el checkpoint base/IFT, compartir los pairs o el reward model, tener un presupuesto de cómputo similar y usar el mismo eval portfolio.
\end{warningbox}

\teachervoice{00:48:45--00:49:35, el expositor usa SteerLM y Starling para señalar que «todavía hay muchos modelos que muestran que PPO/RL supera a DPO». Su conclusión es que la evidencia es mixta, no tomar partido por un bando.}

\subsection{Resumen de la sección}

Zephyr y Tulu 2 demostraron la utilidad práctica de DPO y su escalamiento, mientras que SteerLM/Starling mantuvieron la competitividad de PPO. En la alignment abierta, la unidad experimental adecuada no es el nombre del algoritmo, sino la recipe completa. La siguiente sección entra en el ecosistema de 2024 y analiza la llegada de más participantes, Llama 3, la brecha entre modelos abiertos y cerrados, y el verdadero cuello de botella: la preference data de alta calidad.
